\section{Introduction}
\label{sec:introduction}

The vortex panel method is a numerical technique used in aerodynamics to approximate the potential flow around lifting
bodies of arbitrary shape.
Like the source panel method, it belongs to the family of boundary element methods, where the surface of a body is
discretized into a finite number of panels.
However, each panel is instead assigned a vortex distribution, whose strength is determined such that the flow satisfies
the impermeability boundary condition on the body surface.
The method is based on the assumptions of incompressible, irrotational and inviscid flow, allowing the velocity field to
be represented by a scalar potential that satisfies Laplace's equation.
By superimposing the effects of the vortex panels and the free-stream flow, the velocity potential around the body can
be computed.
Because vortex distributions induce circulation about the body, the vortex panel method is capable of representing
lifting flows, unlike the purely non-lifting source panel formulation.
To yield a physically realistic solution, the panel strengths are further constrained by the Kutta condition, which
enforces smooth flow leaving the trailing edge and ensures a unique circulation is selected.
Once the velocity distribution is known, aerodynamic quantities such as the pressure coefficient and lift can be
obtained using Bernoulli's equation and the Kutta-Joukowski theorem, respectively.
The vortex panel method is particularly useful for predicting pressure distributions and lift on airfoils and other
lifting bodies, and forms the basis for more advanced panel methods that combine sources and vortices to model both
thickness and lifting effects simultaneously.
However, the vortex panel method is very senisitive to the placmenet and
number of panels and can produce incorrect oscillatory results.
